\pdfbookmark[1]{Das bei den Beobachtungsterminen anzuwendende Verfahren.}{verfahren547}
\begin{center}
{\large\itshape Das bei den Beobachtungsterminen anzuwendende Verfahren.\par}
\vspace{5mm}\hrule\vspace{2mm}
{\small Resultate aus den Beobachtungen des magnetischen Vereins. 1836. II.\par}
\vspace{2mm}\hrule\vspace{5mm}
\end{center}
Die sechs jährlich festgelegten Termine fallen gegen Ende der Monate Januar, März, Mai, Juli, September und November. Sie beginnen am letzten Samstag jedes dieser Monate mittags nach Göttinger mittlerer Zeit und enden am Mittag des folgenden Tages. Die bisher an jeden Haupttermin angeschlossenen Nebentermine (am Dienstag und Mittwoch der folgenden Woche jeweils abends von 8 bis 10 Uhr) entfallen künftig.

Bei jedem Termin wird die Stellung der Magnetnadel in der Regel alle fünf Minuten bestimmt, sodass ein Termin 289 Ergebnisse liefert. In Göttingen wird die Uhr vor Beginn jedes Termins genau auf mittlere Zeit gestellt. Da es sehr wünschenswert ist, die einzelnen Bestimmungen an den verschiedenen Beobachtungsorten möglichst gleichzeitig vorzunehmen, stellen die Beobachter an den meisten anderen Orten ihre Uhren ebenfalls gewöhnlich auf Göttinger mittlere Zeit. Wo das nicht ohne Weiteres möglich ist, empfiehlt es sich, zu denjenigen vollen Minuten der Uhr zu beobachten, die den Göttinger Beobachtungszeiten am nächsten kommen. Hätte man etwa vor
\SourceLeaf{548}Beginn des Termins festgestellt, dass die für die Beobachtung verwendete Uhr der Göttinger mittleren Zeit um $13'\,48''$ vorausgeht, so wäre die Stellung der Nadel zu den Uhrzeiten $0^{\mathrm h}\,14'\,.....\,0^{\mathrm h}\,19'\,.....\,0^{\mathrm h}\,24'\,.....\,0^{\mathrm h}\,29'$ und so weiter zu bestimmen. Es empfiehlt sich in jedem Fall, \textit{volle} Minuten zu wählen, weil sich die einzelnen Handgriffe so leichter mechanisch ausführen lassen.

Mit der Stellung der Magnetnadel, die für die einzelnen Zeitpunkte bestimmt werden soll, ist hier diejenige gemeint, die der aufgehängte Magnetstab einnehmen würde, wenn er (oder genauer gesagt seine magnetische Achse) in diesem Augenblick genau im magnetischen Meridian läge. Davon ist die Stellung zu unterscheiden, die er im betreffenden Augenblick tatsächlich einnimmt. Diese Unterscheidung war unnötig, solange man nur Nadeln verwendete, mit denen keine sehr hohe Genauigkeit zu erreichen war. Man musste nur dafür sorgen, dass die Nadel um die Zeit der Beobachtung keine erkennbare Schwingung ausführte, und erhielt damit unmittelbar das Gesuchte. Bei den viel höheren Anforderungen, die man mit den jetzt eingeführten Geräten an die Genauigkeit der Bestimmungen stellen kann und muss, kommt eine solche unmittelbare Bestimmung nicht mehr in Betracht. Es liegt nicht in unserer Macht, die Nadel des Magnetometers so vollkommen zur Ruhe zu bringen, dass keine erkennbaren Schwingungsbewegungen zurückbleiben. Zumindest lässt sich das nicht mit Sicherheit ohne Zeitaufwand erreichen und nicht auf Dauer aufrechterhalten. An die Stelle der unmittelbaren Beobachtung müssen daher mittelbare Bestimmungen treten, bei denen die Nadel nicht vollkommen zur Ruhe gebracht werden muss.

Als Erstes bietet sich die Methode an, die Nadel absichtlich im schwingenden Zustand zu beobachten, zwei aufeinanderfolgende äußerste Stellungen (ein Minimum und ein Maximum) auf der Skala aufzuzeichnen und aus beiden den Mittelwert zu bilden. Dieses an sich einwandfreie Verfahren erfordert jedoch eine Abwandlung, wenn die Schwingungen beträchtlich sind, und ist bei kleinen Schwingungen nur unter einer einschränkenden Bedingung zulässig. Im ersten Fall ist die fortschreitende Abnahme des Schwingungsbogens nämlich selbst von einer Schwingung zur nächsten merklich. Die Abweichung vom wirklichen Meridian auf der Seite des Maximums ist daher bereits kleiner, als sie beim vorhergehenden Minimum auf der gegenüberliegenden \SourceLeaf{549}Seite war. Folglich fällt der Mittelwert aus diesem Minimum und dem folgenden Maximum zu klein aus. Aus demselben Grund ergibt der Mittelwert aus diesem Maximum und dem folgenden Minimum ein zu großes Ergebnis. Da aber die Abnahme des Schwingungsbogens über einige Schwingungen hinweg nahezu gleichförmig bleibt, kann man den Mittelwert aus zwei solchen Mittelwerten als hinreichend genau und als gültig für den Augenblick der zweiten Elongation ansehen. In einer Formel lässt sich das so ausdrücken. Sind $a,b,c$ die Ablesungen bei drei aufeinanderfolgenden Elongationen (gleichviel, ob die erste und die dritte Minima sind und die zweite ein Maximum oder umgekehrt), so gibt $\tfrac14(a+2b+c)$ die Stellung des magnetischen Meridians im Augenblick der Elongation $b$ an.

Bei kleinen Schwingungen ist dieses Verfahren nur zulässig, wenn sich die Deklination in kurzer Zeit nicht merklich ändert. Dann kann man schon den Mittelwert aus zwei aufeinanderfolgenden Elongationen als gültig für den in der Mitte liegenden Augenblick ansetzen. Im entgegengesetzten Fall, also wenn sich die Deklination rasch und beträchtlich ändert, kann das Verfahren ganz unbrauchbar werden.

Die Methode, die Stellung des magnetischen Meridians aus beobachteten Elongationen zu bestimmen, hat jedoch immer den Nachteil, dass die Zeitpunkte, für die das Ergebnis gilt, von den Zeitpunkten abweichen, für die man die Stellung wissen will, oder nur zufällig mit ihnen zusammenfallen. Mag das auch in den meisten Fällen wenig ins Gewicht fallen, so verdient doch offensichtlich ein anderes Verfahren den Vorzug, das diesen Nachteil vermeidet, bequeme Handhabung, Gleichförmigkeit und alle wünschenswerte Genauigkeit vereint und deshalb von allen Teilnehmern an den Terminbeobachtungen befolgt wird.

Dieses Verfahren beruht auf dem Satz, dass der Mittelwert aus zwei Stellungen der Nadel zu Zeitpunkten, die genau eine Schwingungsdauer auseinanderliegen, mit der Lage des magnetischen Meridians zur Mitte zwischen diesen Zeitpunkten übereinstimmt, in welche Teile der Schwingungsperiode die Zeitpunkte auch fallen mögen. Dieser Satz wäre mathematisch exakt, wenn zum einen keine äußeren Ursachen (wie der Luftwiderstand und dergleichen) den Schwingungsbogen\SourceLeaf{550}. nach und nach verkleinerten und zum anderen die etwaige Änderung der Lage des magnetischen Meridians während dieser kurzen Zwischenzeit als \textit{gleichförmig} betrachtet werden dürfte. Der erste Umstand hat jedoch keinen merklichen Einfluss, wenn man das Verfahren stets nur auf sehr kleine Schwingungsbewegungen anwendet. Was den zweiten betrifft, so sind die Änderungen der Deklination während einer so kurzen Zwischenzeit in der Regel schon an sich kaum merklich. Umso eher ist man berechtigt, wenigstens die Gleichförmigkeit der Änderungen während dieser kurzen Zeit gelten zu lassen\HistoricalNote{550}.

Damit ist die Aufgabe von selbst gelöst. Um die der Deklination zur Zeit $T$ entsprechende Stellung der Nadel zu erfahren, braucht man nur die tatsächlichen Stellungen zu den Zeiten $T-\tfrac12t$ und $T+\tfrac12t$ zu beobachten und daraus den Mittelwert zu bilden, nachdem man ihre Bewegungen nötigenfalls vorher durch geeignete Dämpfungsmittel sehr klein gemacht hat. Dabei bezeichnet $t$ die Schwingungsdauer. Um größere Genauigkeit und Sicherheit zu erzielen, begnügt man sich jedoch nicht damit, sondern nimmt noch einige ähnliche Bestimmungen für ein paar Zeitpunkte kurz vor $T$ und ebenso viele danach vor, stets in gleichen Abständen. Unter dieser Voraussetzung wird der Mittelwert aus allen diesen Ergebnissen, sofern die Änderung der Deklination während dieser Zeit als gleichförmig betrachtet werden darf, das für die Zeit $T$ gültige \textit{Endergebnis} sein und zuverlässiger ausfallen als die einzelne Bestimmung für $T$ selbst.

Wenn das Endergebnis zum Beispiel auf fünf Teilergebnissen beruhen soll, zeichnet man am einfachsten die tatsächliche Stellung der Nadel zu den sechs Zeiten
\[
T-\tfrac52t,\quad T-\tfrac32t,\quad T-\tfrac12t,\quad T+\tfrac12t,\quad T+\tfrac32t,\quad T+\tfrac52t
\]
auf. Sind die aufgezeichneten Zahlen $a,b,c,d,e,f$, so ist $\tfrac12(a+b)$ das für die Zeit $T-2t$ gültige Ergebnis. Ebenso sind
\[
\tfrac12(b+c),\quad\tfrac12(c+d),\quad\tfrac12(d+e),\quad\tfrac12(e+f)
\]
die Ergebnisse für die Zeiten $T-t$, $T$, $T+t$, $T+2t$. Der Mittelwert aus diesen Teilergebnissen\SourceLeaf{551}, also der fünfte Teil ihrer Summe, ist als berichtigtes Endergebnis für die Zeit $T$ anzunehmen.

\Needspace{15\baselineskip}
Als Beispiel seien hier die Einzelheiten der Beobachtung in Göttingen vom 17. August 1836 für $15^{\mathrm h}\,30'$ angeführt. Beobachter war Herr Dr.~\textsc{Wappäus}. Für $t$ wurde $20''$ angenommen.
\[
\begin{array}{rrr|r|}
15^{\mathrm h}&29'&10''&865.2\\
&&30&867.5\\
&&50&866.2\\
&30&10&868.0\\
&&30&867.3\\
&&50&868.5
\end{array}
\quad
\left.\begin{array}{r}866.35\\866.85\\867.10\\867.65\\867.90\end{array}\right\}867.16
\]
Die erste Spalte enthält hier die Beobachtungszeiten, die zweite die aufgezeichneten Skalenteile und die dritte den Mittelwert aus je zwei aufeinanderfolgenden Aufzeichnungen, also die für
\[
15^{\mathrm h}\,29'\,20'',\quad15^{\mathrm h}\,29'\,40'',\quad15^{\mathrm h}\,30'\,0'',\quad15^{\mathrm h}\,30'\,20'',\quad15^{\mathrm h}\,30'\,40''
\]
gültigen Teilergebnisse. Daneben steht das für $15^{\mathrm h}\,30'\,0''$ gültige Endergebnis. In diesem Beispiel ist die fortwährende Änderung der Deklination im Laufe der Beobachtungen offensichtlich und wird auch durch die vorhergehenden und folgenden Ergebnisse bestätigt. Die Ergebnisse lauteten
\[
\begin{array}{lrrll}
\text{für}&15^{\mathrm h}&25'\,0''&.\ .\ .\ .\ .\ .&862.82\\
&&35\quad0&.\ .\ .\ .\ .\ .&872.32
\end{array}
\]
Häufiger als so beträchtliche Änderungen ist jedoch eine während eines Beobachtungssatzes nahezu unveränderte Deklination. In solchen Fällen dienen die kleineren oder größeren Schwankungen der Teilergebnisse als Maß für die mehr oder weniger hohe Zuverlässigkeit der Beobachtungen selbst, mag sie nun vom Geschick und von der Aufmerksamkeit des Beobachters, von der Güte des Geräts selbst oder von den mehr oder weniger günstigen äußeren Umständen abhängen.

Das beschriebene Verfahren wird von den meisten Teilnehmern an den Terminbeobachtungen angewandt. Es setzt voraus, dass die Schwingungsdauer der Nadel bekannt ist. Diese hängt bekanntlich sowohl von der Stärke der
\SourceLeaf{552}Magnetisierung der Nadel als auch von der Intensität der horizontalen Komponente der erdmagnetischen Kraft ab und ist daher streng genommen zu verschiedenen Zeiten nicht ganz dieselbe. Eine Anleitung zur genauen Bestimmung der Schwingungsdauer wird später gegeben [S.\ 374 dieses Bandes]. Für den vorliegenden Zweck muss man sie aber nicht sehr genau kennen. Man kann daher die kleinen Änderungen, denen sie unterliegt, außer Acht lassen und darf den genauen Wert sogar durch die nächstgelegene volle Sekunde ersetzen. So fallen die Augenblicke, in denen der Beobachter die Stelle im Skalenbild, die unter dem senkrechten Faden des Fernrohrs erscheint, genau erfassen muss, immer auf volle Sekunden. Das ergibt sich von selbst, wenn die der wahren Schwingungsdauer am nächsten liegende ganze Zahl gerade ist. Ist sie aber ungerade, so hat man, um diesen Vorteil nicht zu verlieren. die Wahl zwischen folgenden drei Mitteln.

I. Man hält sich dennoch an die nächstgelegene gerade Zahl. Das ist umso eher zulässig, je weniger ihr Abstand vom wahren Wert eine halbe Einheit überschreitet, je größer die Schwingungsdauer überhaupt ist und je besser man die Nadel stets in nahezu ruhigem Zustand halten kann. Die Nadel im magnetischen Observatorium zu Göttingen hat zum Beispiel gegenwärtig eine Schwingungsdauer von $20''64$. Obwohl hier 21 die nächstgelegene Zahl ist, kann man sich unter den hier herrschenden Bedingungen, unter denen der Schwingungsbogen selten ein paar Skalenteile überschreitet, meist ohne Bedenken an die bequemere Zahl 20 halten. Denn wie sich leicht zeigen lässt, kann der \textit{daraus} entstehende Fehler eines Teilergebnisses den zwanzigsten Teil des Schwingungsbogens und der Fehler des Endergebnisses den hundertsten Teil nicht überschreiten. Einem Beobachter, dessen Nadel eine Schwingungsdauer von $10''64$ hätte, wäre dagegen zu empfehlen, die Zahl 11 und eine der folgenden Abwandlungen zu wählen, besonders wenn er die Nadel nicht ebenso gut zur Ruhe bringen könnte.

II. Man wählt zwar die ungerade Zahl, legt aber die Beobachtungszeitpunkte, die nach der obigen Formel auf halbe Sekunden fallen würden, entweder alle eine halbe Sekunde später oder alle eine halbe Sekunde früher. Das macht offenbar nur den Unterschied, dass sämtliche Endergebnisse nun für eine halbe Sekunde nach beziehungsweise vor der vollen Minute der Uhrzeit gelten.

\SourceLeaf{553}III. Man stützt das Endergebnis auf eine gerade Anzahl von Teilergebnissen, während das oben entwickelte Verfahren von einer ungeraden Anzahl ausgeht. Dann fallen die Beobachtungszeiten von selbst auf volle Sekunden, ob die nächstgelegene ganze Zahl, die man anstelle der wahren Schwingungsdauer annimmt, nun gerade oder ungerade ist. Soll das Endergebnis zum Beispiel von sechs Teilergebnissen abhängen, so sind die Beobachtungszeiten
\[
T-3t,\quad T-2t,\quad T-t,\quad T,\quad T+t,\quad T+2t,\quad T+3t
\]

Bei diesem Verfahren wird der Einfluss des von der Schwingungsdauer weggelassenen Bruchteils auf das Endergebnis noch vollständiger ausgeschaltet als bei dem zuvor beschriebenen Verfahren. Es ist vor allem Beobachtern zu empfehlen, die kleinere Geräte oder Nadeln mit vergleichsweise kurzer Schwingungsdauer verwenden.

Es sei noch bemerkt, dass das Auflegen eines kleinen Gewichts die Schwingungsdauer der Nadel vergrößert. Wählt man Gewicht und Auflagestelle passend, so lässt sich die Schwingungsdauer einer ganzen Zahl von Sekunden äußerst nahe bringen. Diesen Ausweg haben wohl einige Beobachter gewählt, die es nicht hinreichend in der Hand hatten. etwas größere Schwingungsbewegungen von ihrer Nadel fernzuhalten. Das bleibt jedoch ein sehr unzureichender Notbehelf. Denn auch wenn der obige Satz unter solchen Umständen als ganz exakt gilt, sind die Ergebnisse immer erheblich weniger genau. Wenn sich die Nadel augenfällig bewegt, kann man den einem bestimmten Sekundenschlag entsprechenden Skalenteil und dessen Bruchteil unmöglich ebenso genau erfassen wie bei einer so langsamen Bewegung, dass man eine Änderung in einer Sekunde kaum bemerkt. Die Vorschrift, die Nadel stets in hinreichend ruhigem Zustand zu halten, kann daher gar nicht nachdrücklich genug betont werden.

Eben deshalb ist es wichtig, zwischen zwei aufeinanderfolgenden Beobachtungssätzen stets genügend Zeit zu lassen, um die Nadel nötigenfalls zur Ruhe zu bringen. Bei der Nadel des Göttinger magnetischen Observatoriums beträgt diese Zwischenzeit mit der ersten Methode $3'\,20''$, mit der zweiten würde sie $2'\,54''$ betragen. In beiden Fällen reicht sie für Geübte zum genannten Zweck aus. Gewöhnlich nutzen die Beobachter die Zwischenzeit (da
\SourceLeaf{554}die Nadel nur äußerst selten beruhigt werden muss), um eine Reinschrift der Beobachtung anzufertigen und das Endergebnis zu berechnen. Wo die Nadel dagegen eine viel längere Schwingungsdauer hat und damit die Zwischenzeit zwischen zwei Beobachtungssätzen viel kürzer ausfällt, ist eine Abwandlung der obigen Methoden vorzuziehen, die diesen Nachteil beseitigt.

Die Abwandlung besteht darin, dass man zwischen den einzelnen Beobachtungszeitpunkten jeweils einen von mehreren gleich großen Teilen der Schwingungsdauer (die Hälfte oder ein Drittel) verstreichen lässt, statt die ganze Schwingungsdauer abzuwarten. So kann man die Aufzeichnungen für jeden Satz in kürzerer Zeit erledigen und gewinnt mehr Zeit zwischen zwei Sätzen. Zugleich vermeidet man den Nachteil des früheren Verfahrens, den größten Teil der Zeit zwischen zwei Aufzeichnungen untätig zu sein. Geübtere Beobachter wenden das abgeänderte Verfahren daher auch dort gern an, wo die Schwingungsdauer nicht besonders lang ist. An unserem Observatorium zeichnen mehrere Beobachter in Abständen von $10''$ (als Hälfte von $20''$) oder sogar von $7''$ (als Drittel von $21''$) auf. Einige Beispiele erläutern am besten, was dabei noch zu beachten ist.

\Needspace{15\baselineskip}
\begin{center}\textit{Beobachtung am} 17. \textit{August} 1836, \textit{für} $10^{\mathrm h}\,20'$ \textit{durch Herrn Prof. Ulrich.}\end{center}
\[
\begin{array}{rrr|r|}
10^{\mathrm h}&19'&30''&869.9\\
&&40&871.3\\
&&50&871.7\\
&20&0&870.8\\
&&10&871.0\\
&&20&872.4\\
&&30&872.9
\end{array}\quad
\left.\begin{array}{r}870.80\\871.05\\871.35\\871.60\\871.95\end{array}\right\}871.35
\]
Die zweite Spalte enthält die einzelnen Aufzeichnungen und die dritte die Teilergebnisse. Dabei ist $870.80$ der Mittelwert der ersten und dritten Aufzeichnung und gilt somit für $10^{\mathrm h}\,19'\,40''$ und so weiter. In diesem Beispiel aus einer Zeit rascher Deklinationsänderungen sieht man mit Vergnügen, wie sicher ein geübter Beobachter die Änderungen in 10 Sekunden erkennen kann.

\SourceLeaf{555}\Needspace{16\baselineskip}
\begin{center}\textit{Beobachtung am} 25. \textit{März} 1837, \textit{für} $0^{\mathrm h}\,5'$ \textit{durch Herrn Dr. Goldschmidt.}\end{center}
\[
\begin{array}{rrr|r|}
0^{\mathrm h}&4'&32''&847.3\\
&&39&847.2\\
&&46&847.8\\
&&53&848.7\\
&5&0&848.9\\
&&7&848.1\\
&&14&847.0\\
&&21&846.9\\
&&28&847.3
\end{array}\quad
\left.\begin{array}{r}848.00\\848.05\\847.95\\847.85\\847.90\\847.70\end{array}\right\}847.91
\]
Das erste Teilergebnis entsteht hier durch Zusammenfassen der ersten und vierten Aufzeichnung, das zweite durch Zusammenfassen der zweiten und fünften und so weiter.

In diesen Beispielen war der gewählte Bruchteil der zugrunde gelegten angenäherten Schwingungsdauer eine ganze Zahl. Wo das nicht der Fall ist, muss man die Schwingungsdauer in ungleiche Teile zerlegen. Das hat jedoch keinen Nachteil, wenn man die Aufzeichnungen nur so anordnet, dass zwischen den jeweils zu verbindenden Aufzeichnungen immer \textit{derselbe} angenäherte Wert der Schwingungsdauer liegt, und wenn man, sofern man es der Mühe wert hält, die zum Endergebnis gehörende Zeit mit ihrem Bruchteil im Protokollauszug vermerkt. Will man beispielsweise bei den Beobachtungen mit dem 25-pfündigen Stab in der Sternwarte, dessen Schwingungsdauer jetzt $43''14$ beträgt, den dafür einzusetzenden Näherungswert $43''$ in vier Teile zerlegen und das Endergebnis auf fünf Teilergebnisse stützen, so stellt man die Beobachtungen nach dem folgenden Schema an
\[
\begin{array}{rrr|}
0^{\mathrm h}&4'&17''\\
&&28\\
&&39\\
&&49\\
&5&0\\
&&11\\
&&22\\
&&32\\
&&43
\end{array}\quad
\left.\begin{array}{rrr}0^{\mathrm h}&4'&38''5\\&&49.5\\&5&0.5\\&&10.5\\&&21.5\end{array}\right\}0^{\mathrm h}\,5'\,0''1
\]
Hier enthält die erste Spalte die Aufzeichnungszeiten und die zweite die Zeiten, für die die Teilergebnisse eigentlich gelten. Dabei ist es natürlich
\SourceLeaf{556}ganz gleichgültig, dass das Endergebnis genau genommen auf $0^{\mathrm h}\,5'\,0''1$ fällt. Soll das Endergebnis auf sechs Teilergebnissen beruhen, so gilt folgendes Schema
\[
\begin{array}{rrr|}
0^{\mathrm h}&4'&12'\\
&&22\\
&&33\\
&&44\\
&&55\\
&5&5\\
&&16\\
&&27\\
&&38\\
&&48
\end{array}\quad
\left.\begin{array}{rrr}0^{\mathrm h}&4'&33''5\\&&43.5\\&&54.5\\&5&5.5\\&&16.5\\&&26.5\end{array}\right\}0^{\mathrm h}\,5'\,0''
\]

Am deutlichsten zeigt sich der Vorteil der abgeänderten Beobachtungsweise, wenn man den Verlauf der magnetischen Deklination in kürzeren Abständen als fünf Minuten verfolgen will. Beim gewöhnlichen Verlauf der Deklinationsänderungen reichen diese Abstände aus. Sie sind aber immer noch zu groß, um den stärkeren und schneller wechselnden Änderungen ganz gerecht zu werden. Da sich kürzere Abstände kaum zur allgemeinen und durchgehenden Regel für die vierundzwanzigstündigen Termine machen ließen, gab gerade diese Überlegung den Anlass zur Festsetzung der Nebentermine, bei denen jeweils zwei Stunden lang alle drei Minuten beobachtet werden sollte. An manchen Orten war es jedoch schwierig, diese Nebentermine abzuhalten, und es hat sich auch so ergeben, dass bisher nur bei wenigen davon beträchtliche Bewegungen vorkamen. Daher wurde beschlossen, die Nebentermine von jetzt an entfallen zu lassen, zumal sich derselbe wichtige Zweck auch auf andere Weise und sogar noch besser erreichen lässt. Die Abstände von fünf Minuten bleiben wie bisher die Regel. Sobald man aber eine schnelle Änderung der Deklination bemerkt, folgen die Beobachtungssätze so lange wie nötig im Abstand von $2\frac12$ Minuten aufeinander, jeweils $2\frac12$ Minuten nach dem vorigen. Nach dem oben Erläuterten genügt es, als weitere Veranschaulichung dem obigen Beispiel vom 17. August [S.\ 548] noch die unmittelbar darauf folgende Beobachtung beizufügen.
\SourceLeaf{557}
\[
\begin{array}{rrr|r|}
10^{\mathrm h}&22'&0''&875.0\\
&&10&874.8\\
&&20&876.0\\
&&30&877.1\\
&&40&876.8\\
&&50&876.1\\
&23&0&877.1
\end{array}\quad
\left.\begin{array}{r}875.50\\875.95\\876.40\\876.60\\876.90\end{array}\right\}876.27\quad\text{für }10^{\mathrm h}\,22'\,30''
\]

Alle auswärtigen Teilnehmer werden gebeten, in solchen Fällen ebenso zu verfahren. Es besteht kein Zweifel, dass dann bei allen größeren Bewegungen stets zahlreiche Beobachtungen, die einander bis in die Einzelheiten entsprechen zusammenkommen und interessante Aufschlüsse über die Verhältnisse dieser bemerkenswerten Erscheinungen geben werden.

Wer bei der Beobachtung eine Uhr verwendet, die andere Zeitabschnitte schlägt als eine Sekundenpendeluhr, wird keine besondere Anleitung brauchen. Man zählt dann anstelle der Sekunden die Schläge der Uhr und richtet die Arbeit ganz entsprechend so ein, dass alle Beobachtungen bei bestimmten Schlägen vorgenommen werden. Die Schläge eines Chronometers stets richtig zu zählen, erfordert jedoch etwas mehr Aufmerksamkeit als bei einer Pendeluhr, besonders wenn der Zeiger des Chronometers etwas exzentrisch sitzt und deshalb nicht überall auf dem Zifferblatt, wo er es sollte, genau auf die Sekundenstriche springt.

Einige allgemeine Vorsichtsmaßnahmen verdienen hier noch Erwähnung, auch wenn sie zum Teil scheinbare Kleinigkeiten betreffen. Mancher angehende Beobachter könnte sie anfangs leicht übersehen, wenn man ihn nicht vorher darauf aufmerksam macht.

Die allererste Voraussetzung ist, dass die Nadel völlig frei schwingen kann. Hindernisse für die freie Bewegung, die sofort ins Auge fallen, wird natürlich jeder Beobachter von selbst zu beseitigen wissen. Es gibt aber auch andere, die dem Auge fast entgehen und die Beobachtungen trotzdem ganz verderben können.

In der wärmeren Jahreszeit gelangt zuweilen eine Spinne in den Kasten (am leichtesten, wenn die Seitenöffnung vor dem Spiegel stets offen bleibt). Sie spinnt ein Netz oder einen einzelnen Faden zwischen dem Magnetstab oder seinem Zubehör und dem Kasten und hemmt dadurch die freie Bewegung des
\SourceLeaf{558}Magnetstabs. Man tut daher gut daran, sich kurz vor jedem Termin zu vergewissern, dass der Kasten innen sauber ist. Hat der Kasten einen verglasten Deckel, so erkennt man größere Insekten oder Netze schon von außen. Man sollte es dennoch nicht versäumen, den Deckel abzuheben und genauer nachzusehen. Selbst wenn man gar keinen Faden sieht, sollte man sich nicht darauf verlassen. Denn wie wiederholte Erfahrungen gezeigt haben, reicht schon der feinste Faden, der mit bloßem Auge unsichtbar oder nur bei ganz besonderer Beleuchtung erkennbar ist, aus, um die freie Bewegung zu hemmen und die Beobachtungen zu verderben. Um sich vor diesem gerade wegen seiner Unsichtbarkeit besonders gefährlichen Feind zu schützen, fährt man mit dem Finger, einem Stäbchen oder dergleichen auf allen Seiten am Magnetstab entlang, rechts, links, vorn, hinten, oben und unten. Dadurch wird ein etwa vorhandener Faden zerrissen. Fast ebenso sicher erreicht man dieselbe Wirkung, indem man den Stab in sehr große Schwingungen versetzt. Es verdient noch erwähnt zu werden, dass solche und ähnliche Hindernisse für die freie Bewegung stets mit einer Verkürzung der Schwingungsdauer der Nadel verbunden sind. Schon äußerst zarte Spinnenfäden bewirken eine sehr beträchtliche Verkürzung der Schwingungsdauer, wofür weiter unten ein bemerkenswertes Beispiel folgen wird.

Für die nächtlichen Beobachtungen muss man die Skala beleuchten. Bei den Terminbeobachtungen in Göttingen geschieht das mit zwei \textsc{Argand}-Lampen. Über der Flamme steigt stets erwärmte Luft auf. Steht eine der Lampen fast genau unter dem Fernrohr, beeinträchtigt diese Luftströmung vor dem Objektiv die Sicht. Die Teilstriche der Skala scheinen zu zittern oder zu wallen. Dieser Übelstand trat bei den ersten Beobachtungen in Göttingen öfter auf und ist ganz verschwunden, seit jede Lampe einen seitwärts gebogenen Schornstein aus Kupferblech hat.

Da sich die Arbeit bei den Terminbeobachtungen stets auf eine größere oder kleinere Zahl von Personen verteilen muss, wird ihre Sehweite gewöhnlich beträchtlich verschieden sein. Völlig deutliches Sehen ist aber für gute Beobachtungen eine unerlässliche Voraussetzung. Wird ein Weitsichtiger, für dessen Auge das Fernrohr auf vollkommen deutliches Sehen eingestellt war, von einem Kurzsichtigen abgelöst, so könnte dieser ohne eine Änderung am Fernrohr keine brauchbaren Beobachtungen
\SourceLeaf{559}anstellen. Eine zusätzliche Konkavlinse wäre unbequem und auch wegen des beträchtlichen Lichtverlusts nicht zu empfehlen. Es reicht nicht aus, nur das Okularrohr hineinzuschieben. Dadurch würde zwar das Skalenbild deutlich, doch das Fadenkreuz bliebe unscharf und hätte gegenüber dem Bild des Gegenstands eine Parallaxe. Bei der Bauweise der gewöhnlich für solche Beobachtungen verwendeten Fernrohre müsste man deshalb zugleich die innere Hülse, die das Fadenkreuz trägt, im Okularrohr verschieben und näher an das Okularglas heranbringen. Das erfordert aber eine geübte Hand, kostet Zeit und ist auch aus anderen Gründen im vorliegenden Fall nicht zu empfehlen. Man kann dem Bedürfnis aber sehr einfach abhelfen, indem man sich folgendes Verfahren zur Regel macht. Vor Beginn der Beobachtungen stellt man das Okularrohr im Fernrohr und das darin befindliche Fadenkreuz so ein, dass der Kurzsichtigste unter den Beobachtern Fadenkreuz und Skalenbild zugleich vollkommen deutlich sieht. Kommt ein weitsichtiger Beobachter an die Reihe, so braucht er das dem Auge nächste Glas nur so weit herauszuschrauben, dass er das Fadenkreuz vollkommen scharf sieht. Das Okularrohr und das Fadenkreuz darin lässt er unverrückt. Damit ist auch das Skalenbild von selbst völlig deutlich zu sehen. Ein Kurzsichtiger, der später an die Reihe kommt, braucht das Glas dann nur so weit wieder hineinzuschrauben, wie sein Auge es erfordert.

Zur Prüfung, ob das Fernrohr unverrückt geblieben ist, dient eine Marke in einer solchen Entfernung, dass sie bei der für ein deutliches Skalenbild erforderlichen Okularstellung ebenfalls deutlich erscheint. Im Göttinger magnetischen Observatorium besteht sie lediglich aus einem feinen senkrechten Strich an der Nordwand\HistoricalNote{559}.

\SourceLeaf{560}Vor Beginn der Beobachtungen richtet man das Fernrohr auf die Marke. Danach wiederholt man die Prüfung von Zeit zu Zeit und bringt die optische Achse des Fernrohrs, sobald sich eine Abweichung zeigt, wieder in die vorige senkrechte Ebene. Hat man zu beiden Seiten der Marke noch eine Einteilung angebracht, so erkennt man daran zugleich die Größe der erforderlichen Korrektur. Dabei ist jedoch zu bedenken, dass diese Teilungen, selbst wenn sie ebenso wie die Skalenteile Millimeter sind, genau genommen in Bogensekunden einen etwas anderen Wert haben als die Skalenteile. Fehlt diese Hilfseinteilung, so kann man sich aber damit begnügen, die gefundene Abweichung nach Augenmaß in Skalenteilen zu schätzen, so wie diese erscheinen.

Die Beobachtungen werden am senkrechten Faden des Fadenkreuzes vorgenommen. Der waagerechte Faden dient lediglich dazu, ungefähr die Mitte des senkrechten Fadens zu bezeichnen. Damit es keinen Unterschied macht, ob man die Skalenteile etwas höher oder tiefer im Gesichtsfeld erscheinen lässt, muss das Fadenkreuz so eingestellt sein, dass das Bild eines feststehenden Gegenstands am Schnittpunkt der Fäden genau auf dem senkrechten Faden bleibt, wenn man das Fernrohr etwas auf und nieder bewegt. Auch zu dieser Berichtigung dient die Marke. Lässt man die Stellung des Okularrohrs unverändert, braucht man die Berichtigung übrigens nur selten zu wiederholen.

Der von der Mitte des Objektivs herabhängende Lotfaden liegt so nahe an der Skala, dass die Bilder von beiden im Fernrohr gleich deutlich erscheinen. Daher kann man den Teilstrich, den der Faden verdeckt, sehr genau beobachten. Man bringt die Skala so an, dass dieser Punkt ihre Mitte ist oder ein dafür frei gewählter Teilstrich. Während der Beobachtungen muss man von Zeit zu Zeit prüfen, ob die Skala unverrückt geblieben ist. Findet man eine kleine Änderung, so braucht man die Skala aber nicht in ihre vorige Stellung zurückzubringen. Es genügt, den dem Lotfaden entsprechenden Teilungspunkt im Protokoll zu vermerken.

\SourceLeaf{561}Dabei ist es vielleicht sinnvoll, auf ein paar Kleinigkeiten besonders aufmerksam zu machen.

Es wird zwar vorausgesetzt, dass Magnetometer und Fernrohr so aufgestellt sind, dass der mittlere Stand der magnetischen Deklination ungefähr der Mitte der Skala entspricht. Bei beträchtlichen Schwankungen gerät diese Mitte jedoch nicht selten ganz aus dem Gesichtsfeld, sodass man die oben beschriebene Prüfung nicht vornehmen kann. Hat man in einem solchen Augenblick Anlass zu dieser Prüfung, so muss man den Beruhigungsmagneten zu einem Zweck einsetzen, der seinem üblichen Gebrauch gerade entgegengesetzt ist. Man versetzt die Magnetometernadel in Schwingungen, die bis zur gesuchten Stelle oder ein wenig darüber hinaus reichen. So kann man den Lotfaden in der Mitte des Gesichtsfeldes sehen, und zwar in einem Abschnitt der Schwingungsperiode, in dem die Bewegung langsam ist und das genaue Erfassen des entsprechenden Teilungspunkts daher nicht behindert. Tritt das während der Beobachtungen ein, muss man die Nadel sogleich wieder zur Ruhe bringen, um den folgenden Beobachtungssatz möglichst nicht zu versäumen. Daran zeigt sich, wie nützlich es ist, mit dem Gebrauch des Beruhigungsmagneten gründlich vertraut zu sein.

Im umgekehrten Fall, also immer wenn die Deklination in die Nähe der Skalenmitte fällt, ist für Ungeübte eine andere Warnung nötig. Sie dürfen den Lotfaden nicht mit dem senkrechten Faden des Fernrohrs verwechseln. An unserem Gerät sehen sich beide tatsächlich so ähnlich, dass man sie bei sehr ruhiger Nadel ohne ein paar am letztgenannten Faden befestigte Stäbchen wohl verwechseln könnte. An einem anderen Ort ist es früher tatsächlich einmal vorgekommen, dass ein Beobachter die Nadel eine halbe Stunde lang für völlig unbeweglich hielt, während er die ganze Zeit den falschen Faden beobachtet hatte. Weil es die Beobachtung stets ein wenig erschwert, wenn die Fäden sehr nahe beieinanderliegen, entfernt man in einem solchen Fall den Lotfaden am besten für eine Weile.

Was die Form der Mitteilung betrifft, so pflegen einige die Beobachtungen ganz \textit{in extenso} einzusenden, andere die Teil- und Endergebnisse und mehrere nur die Endergebnisse. Sofern man zuvor die Rechnungen durchgesehen und die mitgeteilten Zahlen mit den Originalen abgeglichen hat, kann dieser Auszug genügen. Die Beobachtungen selbst müssen jedoch aufbewahrt werden, um erforderlichenfalls
\SourceLeaf{562}darauf zurückgreifen zu. können. Für Zeiten, in denen ungewöhnlich starke Bewegungen vorkommen, bleibt die sofortige vollständige Mitteilung jedoch wünschenswert. Außer den Beobachtungszahlen sind die sonstigen damit verbundenen Umstände, der Wert der Skalenteile (oder die Messungen, auf denen die Bestimmung beruht), die Schwingungsdauer, Stand und Gang der Uhr, die Namen der Beobachter, Erläuterungen zu etwa als zweifelhaft bezeichneten Beobachtungen und dergleichen beizufügen. Dass schließlich stets eine \textit{baldige} Einsendung erwünscht sein muss, braucht nicht eigens erwähnt zu werden.

